\section{Experimental design}
\label{sec:experimental-design}

The experiments measure two things about the protocol of Section~\ref{sec:proposed-method}: what each way of breaking it costs in the validation estimate, and what each gate sees. Every comparison is paired: the protocol and each leaked variant are fitted on the same data draw with the same seed, so that a difference between them is the effect of the row set that one step read and nothing else.

\subsection{Predictions}

Nine predictions are fixed, each with a decision rule in terms of a measured quantity applied by the table script to the cells of the paper that the prediction names (the caption of Table~\ref{tab:hypotheses} lists them; H1, H2 and H4 include the product's own scorecard, which the paper reports with the protocol, with the fit on every row and with the calibration map), and a prediction that fails is reported as not held. They were not stated before the results and were not registered outside the repository: the author wrote them in the table script between 4 and 7 November 2026, during the analysis (the record does not date the writing more closely), together with the experiment outputs that decide them, and the outputs were committed with them. For two of them the record shows that part of the answer was already known. The committed document of the case C22 (31 October 2026) held an AUC-test power of \nv{c22.auc.0.5}\% at $s=0.5$ and a stability-index power of \nv{c22.csi.1}\% at $s=1$, and H6 asks for at least 97\% and at most 20\%; the committed documents of the case C03 (30 October 2026) held out-of-time gaps of the scorecard whose smallest absolute value is \nv{c03.mingap} points, and H8 asks for at least 2. H6 and H8 are therefore checks that new runs reproduce committed values against thresholds placed beside them, not tests of an unseen outcome. For the other seven the record shows no earlier statement, and none is claimed. Table~\ref{tab:hypotheses} gives them with the quantity that decides each and its outcome, which the table script computes from the committed outputs of the experiments by the rule written in the prediction; no outcome is entered by hand. The predictions cover the unbiasedness of the protocol (H1), the direction and size of the bias of the leaks that read holdout labels (H2), the steps that read only holdout inputs (H3), the calibration map (H4), the detection by the gates (H5), the leak that exists only in development (H6), the closed forms of the theory (H7), the random split against the out-of-time cohort (H8) and cross-validation (H9). The thresholds (0.005 for a negligible bias; 97\%, 20\% and 2\% for the power of the three detectors of the planted leak; 1 and 2 percentage points for the level gap) are the decision rules of the design, round values chosen by the author; no proposition predicts the 97\%, 20\% or 2\% levels. The table gives the measured quantity next to each outcome so that a reader can apply other thresholds.

\begin{table*}[htbp]
\centering
\small
\caption{The nine predictions of the experimental design, each with the quantity that decides it and its outcome, computed by the table script from the committed experiment outputs. A prediction that fails is reported as not held.}
\label{tab:hypotheses}
\fitwidth{%
\begin{tabular}{lp{7.2cm}p{4.6cm}l}
\toprule
Id & Prediction stated before the experiments & Evidence & Outcome \\
\midrule
H1 & \relax The proposed protocol's holdout AUC is unbiased for the fitted model's true AUC (the 95\% interval of the bias contains 0) and its 95\% interval covers (the Wilson interval of the coverage contains 0.95); refuted if either fails in any cell. & 6 of 6 cells & held \\
H2 & \relax A step that reads holdout labels (WoE, screen, or the whole fit) inflates the AUC (the 95\% interval of the bias excludes 0 from above); refuted if any of the three does not in a data set. & syn: 3 of 3, small: 3 of 3, tw: 2 of 3 & not held \\
H3 & A step that reads only holdout inputs (imputation, scaling, bin edges) leaves the estimate within 0.005 of the truth. & 12 of 12 cells & held \\
H4 & A calibration map fitted on the holdout biases the AUC and the Brier score optimistically (positive AUC bias, negative Brier bias, intervals excluding 0). & 6 of 6 cells & held \\
H5 & At least one of the three gates fails on every leak planted in the fitting path, in every draw, and none fails on the proposed protocol. & \relax 15 of 15 leaks, false alarms 0\% & held \\
H6 & \relax A leak that exists only in development is flagged by the AUC-against-development test (power at least 97\% at s = 0.5) and not by the CSI of the leaked feature (power at most 20\% at s = 1) or the single-feature screen (at most 2\% at s = 1). & \relax AUC test power 99.0\% at s = 0.5; CSI power 14.2\% and single-feature flag 0.0\% at s = 1 & held \\
H7 & The closed forms of Proposition 4 reproduce the generator's exact development and decision-time AUC and the simulated means within three standard errors. & largest absolute z 1.8 over seven levels & held \\
H8 & A random split inside the training window reports a level gap below 1 point for every cohort, while the out-of-time gap is at least 2 points. & random gap below 1 point in 5 of 5 cohorts; out-of-time gap at least 2 points in 5 of 5 & held \\
H9 & Preprocessing fitted before the folds inflates the cross-validated AUC; fitted inside the folds it does not. & outside-fold bias positive, inside-fold bias null & held \\
\bottomrule
\end{tabular}%
}
\end{table*}


\subsection{Data and their provenance}

\paragraph{Planted data.} The planted data have a known truth. Each row has $p=40$ independent standard normal inputs, five of which are informative with coefficients $(0.9,-0.7,0.6,0.5,-0.4)$ in a logistic model with intercept $-1.3$, so that the event share is \nv{syn.share}\% and the true probability of every row is known; 5\% of the values in ten of the inputs are set missing completely at random. The development sample has \nv{syn.n} rows (and \nv{small.n} rows in the small-sample variant); the split gives a holdout of \nv{syn.hold} rows. A reference sample of \nv{syn.oracle} rows is drawn from the same model in every repetition, and no step reads it. The reference AUC of a fitted model on this sample has a sampling error of \nv{syn.orcse} (Hanley and McNeil standard error), which is mean-zero across repetitions and does not bias the comparison.

\paragraph{Card holders (C01).} The UCI data set of 30,000 credit-card clients in Taiwan, with the default of the next month as outcome, a share of \nv{c01.rate}\% \citep{yeh2009}; the data are distributed under CC BY 4.0 (\url{https://doi.org/10.24432/C55S3H}). Each repetition draws \nv{tw.oracle} rows as the reference sample and \nv{tw.n} other rows as the development sample, using the \nv{tw.p} numeric inputs of the case. The reference sample has a sampling error of \nv{tw.orcse} in the AUC, larger than that of the planted data, which makes the coverage of the interval on this data set slightly understated and leaves the bias unchanged.

\paragraph{SBA 7(a) loans (C03).} The 7(a) loan files of the U.S. Small Business Administration (\url{https://data.sba.gov/dataset/7a-504-foia}, licence ``U.S. Government Works'', files dated 30 June 2026 in their names), with the 5-year charge-off as outcome, in five cohorts of approval fiscal years whose training windows follow equation~\eqref{eq:window}; Table~\ref{tab:splits} gives the windows. The product's fit of the window and the scorecard and gradient-boosting rungs are the ones of the case, with the seed of the committed documents; the experiment reproduces the committed out-of-time AUC of the scorecard to within \nv{e4.maxrepdiff} for every cohort.

\paragraph{The engine's planted leak (C22).} The binormal generator of the engine \texttt{riskvalidation} 1.1.0 with 150 defaulters and 4,850 survivors in each of the development and the current sample, a genuine score with separation $d=1$, and a leaked feature that is shifted by $s\in\{0,0.25,0.5,1,1.5,2,3\}$ standard deviations for defaulters in development and is noise at decision time. The committed document of the case holds the exact truth of the generator at six of these levels, and the closed forms of Proposition~\ref{prop:leakauc} reproduce it.

\paragraph{The product's scorecard and ladder.} The product's own optbinning scorecard \citep{navas2020} with monotone binning, information-value screen and sign check, fitted on planted data with 12 inputs, and the product's whole ladder of seven rungs, on a planted set of \nv{e2.ladder.n} rows.

\subsection{Methods compared}

The published practices are run, not quoted. Each is a variant of the compact pipeline of Section~\ref{sec:proposed-method} in which one step reads the holdout, so that it is the practice that the cited work describes and nothing else. Imputation and scaling on all development and holdout rows is the pooled preprocessing that the standard textbooks warn against \citep{hastie2009,moscovich2022}. Bin edges on all rows is the same for the supervised binning of a scorecard \citep{siddiqi2017}. The weight of evidence and the association screen computed with the holdout labels are the selection bias of Ambroise and McLachlan \citep{ambroise2002}. Everything fitted on training and holdout is the model whose holdout is its own training data. The configuration chosen on the holdout is the practice of choosing the candidate by the validation AUC and reporting the AUC on the same partition \citep{cawley2010,autobayesian2026}. The calibration map fitted on the holdout, and a calibration slice that overlaps the holdout by 20\% of its rows, are the faulty uses of the held-out data that calibration requires \citep{niculescu2005,zadrozny2002}. Two five-fold procedures complete the set: preprocessing inside each training fold and preprocessing on all development rows before the folds \citep{ambroise2002,varma2006}. The detectors of the planted leak are the engine's own: the ECB discrimination test of the current AUC against the development AUC, with the exact development AUC of the generator passed as the known development value, which is the reading of the ECB instructions \citep{ecb2019}, the characteristic stability index of the leaked feature \citep{yurdakul2020}, the single-feature AUC screen with its threshold of 0.95, and the resemblance AUC between the development and the current feature, the engine's implementation of the idea of the sample-similarity tests of the open toolkits \citep{mercury2023,probatus2026}. The toolkits themselves are not run (Section~\ref{sec:limitations}).

\subsection{Metrics}

The bias of the holdout AUC, the Brier score and the expected calibration error is the mean over repetitions $r=1,\dots,R$ of the estimate on $H$ minus the value of the same fitted model on the reference sample,
\begin{equation}
\hat b=\frac1R\sum_{r=1}^R\Bigl(\hat A_H^{(r)}-\hat A_{\mathrm{ref}}^{(r)}\Bigr),\qquad
\mathrm{se}(\hat b)=\frac{\mathrm{sd}_r\bigl(\hat A_H^{(r)}-\hat A_{\mathrm{ref}}^{(r)}\bigr)}{\sqrt R},
\label{eq:biashat}
\end{equation}
reported as $\hat b\pm1.96\,\mathrm{se}(\hat b)$, with the root mean squared error $\bigl(\frac1R\sum_r(\hat A_H^{(r)}-\hat A_{\mathrm{ref}}^{(r)})^2\bigr)^{1/2}$. The Brier score is $\frac1{|H|}\sum_{i\in H}(\hat p_i-y_i)^2$ and the calibration error is
\begin{equation}
\mathrm{ECE}=\sum_{g=1}^{10}\frac{|H_g|}{|H|}\,\bigl|\bar p_g-\bar y_g\bigr|,
\label{eq:ece}
\end{equation}
over ten groups of equal size of the predicted probability. The coverage of the engine's interval is the share of repetitions with $|\hat A_H-\hat A_{\mathrm{ref}}|\le1.96\,s$, with its Wilson interval \citep{brown2001}. The detection rate of a gate or a test is the share of draws in which it rejects, with a Wilson interval. For the out-of-time comparison the gap is the observed minus the predicted 5-year rate,
\begin{equation}
\mathrm{gap}=\bar y-\bar p,
\label{eq:gap}
\end{equation}
in percentage points, and the AUC is that of the same fitted rungs on the two samples. A closed form of the theory is checked by the standardised difference
\begin{equation}
z=\frac{\hat\mu_{\mathrm{sim}}-\mu_{\mathrm{formula}}}{\mathrm{se}(\hat\mu_{\mathrm{sim}})},
\label{eq:zscore}
\end{equation}
which has to stay within three for the formula to be accepted. These choices follow the practice for simulation studies \citep{morris2019}.

\subsection{Protocol, repetitions and seeds}

Table~\ref{tab:design} lists the experiments, their data, repetitions and seeds. Every repetition has its own seed, derived from the experiment seed and the repetition number by a seed sequence, so that a repetition can be rerun alone and the results do not depend on the number of workers. The statistics are those of the metrics above: 95\% half-widths over repetitions, Wilson intervals for proportions, and the standardised difference for the closed forms. The number of repetitions follows from the resolution that the comparisons need: a half-width of 0.002 in the AUC bias on 2,500 rows separates the leaks that cost more than 0.004 from the protocol, and 100 draws of the gates give a Wilson upper bound of \nv{e2.fa.upper}\% for the false-alarm rate when none occurs. The out-of-time experiment has one fit per cohort, since the longest cohort took \nv{run.e4max} seconds, and its uncertainty is the binomial error of the 5-year rates.

\begin{table}[htbp]
\centering
\small
\caption{The experiments, the data each runs on, the number of repetitions and the seed; every repetition draws its own data, split and reference sample from the seed of the experiment and its own index, and all variants of one repetition share the draw.}
\label{tab:design}
\fitwidth{%
\begin{tabular}{llrr}
\toprule
Experiment & Data and what varies & Repetitions & Seed \\
\midrule
Bias of the estimate & planted, \nv{syn.n} rows, nine leaks and two five-fold procedures & \nv{syn.reps} & \nv{syn.seed} \\
 & planted, \nv{small.n} rows & \nv{small.reps} & \nv{small.seed} \\
 & card holders, \nv{tw.n} rows from 30,000 & \nv{tw.reps} & \nv{tw.seed} \\
 & product scorecard, 12 inputs, two leaks & \nv{prod.reps} & \nv{prod.seed} \\
Gates & planted, \nv{e2.n} rows, every leak, three gates and the control & \nv{e2.reps} & 20261108 \\
 & discrete choice, one to five perturbation draws & \nv{e2.draws.reps} & 20261108 \\
Planted leak & engine generator, seven leak levels, four detectors & \nv{e3.reps} & 20261109 \\
Closed forms & Propositions~\ref{prop:detection}, \ref{prop:selection}, \ref{prop:leakauc} and \ref{prop:psi} & 10,000 to 20,000 & 20261110 \\
Cross-validation & planted, nine combinations of rows and inputs & \nv{e5.reps} & 20261111 \\
Out-of-time & SBA loans, five cohorts, one fit per cohort & 1 & 42 \\
\bottomrule
\end{tabular}%
}
\end{table}

\subsection{Ablations planned}

Two ablations are planned and reported in Section~\ref{sec:ablation}. The first removes one component of the protocol at a time (statistics fitted on the training slice only, configuration chosen inside the training slice, calibration on a slice disjoint from the test slice, preprocessing inside the folds) and reads the cost from the bias and the coverage of the corresponding leak. The second keeps every subset of the three gates and reads, for the leaks planted in the fitting path, the number that at least one gate of the subset detects in every draw, the number that no gate of the subset detects, the mean detection and the false-alarm rate on the protocol. In both ablations the comparison is paired, with the protocol on the same draws as the reference. The first ablation takes, for each rule, the worst leak that its removal allows, so that a rule that is cheap to remove under one leak is not credited with the cost of another; the second reports the share of draws in which a subset detects and not a single threshold, so that a reader can set the level of detection that a pipeline requires.
